\section{Discussion}

\subsection{Limitations}

The results reported in this paper are subject to several limitations that must be acknowledged. First, even the Tatoeba corpus is substantially smaller than the corpora used in WMT competitions (approximately 324k training pairs versus millions), which restricts the ceiling of achievable translation quality. The word-level vocabulary of the diagnostic configuration (12,843 English and 23,568 German tokens) suffered from non-negligible out-of-vocabulary rates, particularly for German---a morphologically rich language with productive inflection; the final configuration mitigates this via SentencePiece BPE with an 8k-piece vocabulary per language, but rare-word handling remains bounded by the corpus size. Second, the main configuration uses a model with approximately 13.5 million parameters, which is one order of magnitude smaller than the Transformer-Base configurations used in Vaswani et al. (2017) ($\approx$ 65 M for $d_{\text{model}}=512$, 6 layers). This capacity gap directly limits the model's ability to capture complex syntactic and semantic regularities.

Third, our evaluation is restricted to a single random seed (42). While this is sufficient to demonstrate that the implementation works correctly, a robust experimental evaluation would require reporting mean and standard deviation across at least 3--5 seeds to account for the variance introduced by random initialization and data ordering. Fourth, decoding results are reported for greedy search and beam search with beam size 4 under sacrebleu-default conventions; more advanced decoding techniques (length-penalty tuning, coverage penalties, checkpoint averaging) remain unexplored. Fifth, the comparison against \texttt{nn.Transformer} (Section~6.7) is not fully controlled: the two runs differ in batch size and warmup, so the measured margin mixes implementation and schedule effects. Sixth, the attention interpretability analysis in Section~5.5 is restricted to four short test sentences (six content tokens each) and to the last layer of a 4-layer stack, which limits the generality of the qualitative claims.

\subsection{CPU and GPU Considerations}

The current implementation is agnostic to the compute device and runs on both CPU and GPU. The main run reported in this paper (hand-rolled + BPE, Pre-LN, 25 epochs) was executed on a single Google Colab GPU (NVIDIA T4) in \textbf{$\approx$ 2.5 hours}; the word-level diagnostic completed in $\approx$ 56 minutes for 8 epochs, and the baseline in $\approx$ 2.3 hours. On a comparable consumer-grade GPU, the same configurations would train in a similar wall-clock time; CPU-only execution at this corpus scale would be impractical, with the GPU speedup coming primarily from batched matrix multiplications in the attention and feed-forward layers. The hand-rolled implementation does not currently exploit mixed-precision training (\texttt{torch.cuda.amp}) or gradient checkpointing, both of which are standard techniques for reducing GPU memory usage and training time in larger models. Integrating these optimizations is left to future work.

\subsection{What the Results Demonstrate}

Despite these limitations, the results make a clear statement: a from-scratch, hand-rolled implementation of the Transformer architecture is capable of learning a nontrivial translation task from raw parallel data, reaching BLEU 44.00 greedy / 45.30 with beam search once properly configured and trained (Table~\ref{tab:results}). The qualitative analysis of Section~5.5 confirms that the attention maps behave as expected at the structural level (strictly lower-triangular decoder self-attention; non-trivial diagonal mass in encoder self-attention), providing confidence in the implementation's internal consistency.

\textbf{Diagnosis of the word-level control run: sub-training, not a flawed implementation.} The initial word-level Post-LN run (8 epochs, val\_loss = 5.01, val\_ppl = 149.84, test BLEU = 0.38)---retained in this paper as a \textit{diagnostic control}---was unambiguously under-trained, not architecturally broken. Three pieces of evidence supported this:

\begin{enumerate}
\item \textbf{The validation loss was still decreasing} at the end of the 8-epoch run, with no sign of plateau; the model simply had not run long enough (the training log is released alongside the code).
\item \textbf{The model had not learned to use cross-attention} to condition its output on the source. Section~5.3 shows that this run collapses to three recurring hypotheses (\textit{hast du das ?}, \textit{sie sind sehr gro\ss{} .}, \textit{ich habe mich gesehen .}) on six unrelated inputs, and Section~5.5 shows that its cross-attention maps on test sentences are diffuse rather than peaked---in contrast to the sharp alignment peaks produced by the final checkpoint. A correct but under-trained network exhibits exactly this signature: the decoder falls back to a high-frequency prior because the encoder representations have not yet been shaped into useful conditioning signals.
\item \textbf{The implementation itself is correct.} All 17 unit tests pass, the encoder/decoder mask shapes and values are validated, the greedy-decode loop terminates at \texttt{<eos>}, and the forward pass returns the attention tensors with the documented shapes. A bug of the kind this work sets out to avoid (incorrect masking, broken softmax, off-by-one in positional encoding) would manifest as NaN losses or shape mismatches, not as a slowly-converging model with a stable architecture.
\end{enumerate}

\textbf{Resolution.} The corrections motivated by this diagnosis---pre-layer normalization, a corrected warmup schedule (4{,}000 steps), subword tokenization (SentencePiece BPE, 8k pieces per language), and a longer 25-epoch budget at batch size 64---jointly take the same implementation from BLEU 0.38 to \textbf{44.00 greedy / 45.30 beam=4} (chrF2 62.41/63.37) at validation perplexity 9.29, confirming that the failure was one of configuration and training rather than of architecture or implementation.

The final results fall within the range typically reported for compact subword models on this corpus scale (roughly BLEU 35--45), making them competitive with published results for this regime, although direct comparisons remain approximate due to differences in splits and tokenization. The remaining gap factors are those discussed above: single-seed evaluation, model capacity, and untuned decoding.

\subsection{Future Directions}

Beyond the ablation studies outlined in Section~6, several extensions would strengthen the contribution of this work. \textbf{First}, replicating the main configuration across 3--5 random seeds and reporting mean $\pm$ standard deviation would establish statistical robustness. \textbf{Second}, retraining the \texttt{nn.Transformer} baseline with fully matched batch size and warmup would isolate the implementation-only effect that the current comparison (Section~6.7) leaves confounded with schedule differences. Third, running the remaining controlled ablations of Table~\ref{tab:ablation} (label smoothing, warmup, model depth/width) under a fixed schedule would quantify each factor's standalone contribution. Fourth, extending the attention analysis to longer test sentences, to all layers, and to a full characterisation of all 8 heads would broaden the interpretability findings beyond short inputs. Fifth, integrating pre-trained embeddings (e.g., fastText or multilingual BERT) as initialization could accelerate training and improve quality on the small Tatoeba corpus. Sixth, extending the implementation to additional language pairs (e.g., EN$\rightarrow$ES, EN$\rightarrow$FR, or DE$\rightarrow$FR) would test the generalizability of the approach. Seventh, compiling the model with \texttt{torch.compile} (PyTorch 2.0+) would provide a direct measure of the performance overhead of the hand-rolled implementation versus the compiled library version.

\subsection{Pedagogical Value}

We emphasize that the primary contribution of this work is not a new architecture or a superior translation system, but a transparent, auditable reference implementation. The combination of clean code, comprehensive unit tests, and reproducible experimental logs makes this repository suitable for use as a teaching resource in graduate-level NLP courses, as a correctness baseline for library developers, and as a starting point for researchers who wish to experiment with architectural modifications without navigating opaque library abstractions. The fact that attention weights are exposed as first-class outputs (Section~3.4) means the implementation is also a useful platform for interpretability research at a level of detail that library-based implementations often hide.
